\documentclass[12pt, a4paper]{article}

\usepackage[T1]{fontenc}
\usepackage[utf8]{inputenc}
\usepackage{lmodern}
\usepackage[a4paper, margin=2.5cm]{geometry}
\usepackage{amsmath}
\usepackage{graphicx}
\usepackage{xcolor}
\usepackage{booktabs}
\usepackage{tabularx}
\usepackage{array}
\usepackage{colortbl}
\usepackage{enumitem}
\usepackage{fancyhdr}
\usepackage{titlesec}
\usepackage{abstract}
\usepackage{parskip}
\usepackage{microtype}
\usepackage{hyperref}
\usepackage{caption}
\captionsetup{labelsep=period}
\usepackage{float}

\hypersetup{
    colorlinks=true,
    linkcolor=blue!70!black,
    urlcolor=blue!70!black,
    citecolor=blue!70!black
}

% ── Page style ────────────────────────────────────────────────────────────────
\pagestyle{fancy}
\fancyhf{}
\lhead{\small\textit{ClinGen-MoE — Use Case Study: PhysioNet 2012}}
\rhead{\small\textit{Said Abolhassan Razavi}}
\cfoot{\thepage}
\renewcommand{\headrulewidth}{0.4pt}

% ── Section formatting ────────────────────────────────────────────────────────
\titleformat{\section}{\large\bfseries}{}{0em}{}[\titlerule]
\titleformat{\subsection}{\normalsize\bfseries}{\thesubsection.}{0.5em}{}
\titlespacing{\section}{0pt}{18pt}{8pt}
\titlespacing{\subsection}{0pt}{12pt}{4pt}

% ── Table column types ────────────────────────────────────────────────────────
\newcolumntype{L}[1]{>{\raggedright\arraybackslash}p{#1}}
\newcolumntype{C}[1]{>{\centering\arraybackslash}p{#1}}
\definecolor{headergray}{HTML}{374151}
\definecolor{rowlight}{HTML}{F9FAFB}

% ─────────────────────────────────────────────────────────────────────────────
\begin{document}

% ── Title ─────────────────────────────────────────────────────────────────────
\begin{titlepage}
    \centering
    \vspace*{3cm}

    {\Huge\bfseries Use Case Study\\[10pt]
    Synthetic ICU Patient Data from PhysioNet 2012}

    \vspace{1cm}
    \rule{0.5\linewidth}{0.4pt}
    \vspace{1cm}

    {\large Said Abolhassan Razavi}\\[6pt]
    {\normalsize Université Paris-Saclay}\\
    {\normalsize Master 1 --- Artificial Intelligence}\\[6pt]
    {\normalsize Project: ClinGen-MoE}\\[6pt]
    {\normalsize\textit{PhysioNet Computing in Cardiology Challenge 2012}}

    \vspace{2cm}

    \begin{abstract}
    \noindent
    This document presents three use cases for synthetic ICU patient data generated
    from the PhysioNet 2012 dataset using the ClinGen-MoE framework. Unlike
    standard synthetic data approaches that handle time-series measurements alone,
    PhysioNet 2012 provides two distinct modalities: static demographic descriptors
    (age, gender, ICU type, height, weight) and time-series clinical measurements
    (24 features over 48 hours). The generative model is conditioned on patient
    demographics, producing synthetic trajectories that are consistent with
    the patient's demographic profile. Three use cases are presented that
    exploit this multimodal structure for clinical AI applications:
    rare subgroup augmentation, ICU-type alert calibration, and
    privacy-safe inter-hospital data sharing.
    \end{abstract}

    \vfill
    {\small\today}
\end{titlepage}

% ─────────────────────────────────────────────────────────────────────────────
\section*{Introduction}

The ClinGen-MoE model generates synthetic ICU patient records that are
statistically faithful to the real PhysioNet 2012 dataset. Synthetic means
are within 2\% of real values across all 24 clinical features. The core
contribution exploited in these use cases is the \textbf{conditional generation}
mechanism: the model takes as input a patient's static descriptors (age, gender,
ICU type, height, weight) and generates time-series vital signs and laboratory
measurements that are physiologically consistent with that demographic profile.

This conditioned generation is fundamentally different from unconditional
synthetic data approaches. It enables three classes of clinical AI applications
that are described in the following sections.

\vspace{8pt}

\begin{table}[H]
\centering
\renewcommand{\arraystretch}{1.4}
\begin{tabularx}{\linewidth}{L{3cm} X C{2cm}}
\toprule
\rowcolor{headergray}
\textcolor{white}{\textbf{Modality}} &
\textcolor{white}{\textbf{Features}} &
\textcolor{white}{\textbf{Count}} \\
\midrule
\rowcolor{rowlight}
\textbf{Static} (demographic) &
Age, Gender, Height, Weight, ICU Type
(CCU, CSRU, MICU, SICU) & 5 \\
\textbf{Vital signs} (time-series) &
Heart Rate, Systolic BP, Diastolic BP, Mean BP,
Respiratory Rate, O$_2$ Saturation,
GCS Verbal, GCS Eye, GCS Motor & 9 \\
\rowcolor{rowlight}
\textbf{Lab results} (time-series) &
Glucose, Potassium, Sodium, Creatinine,
Urea Nitrogen, Bicarbonate, Hematocrit,
White Blood Cells, Magnesium, Platelet Count,
and 5 additional markers & 15 \\
\midrule
\textbf{Total} &
All measured over a 48-hour ICU window & \textbf{24 + 5} \\
\bottomrule
\end{tabularx}
\end{table}

% ─────────────────────────────────────────────────────────────────────────────
\section{Use Case 1: Rare Demographic Subgroup Augmentation}

\subsection{Motivation}

Training reliable clinical AI models requires sufficient data in each patient
subgroup of interest. In practice, many ICU subgroups are too rare:
a small hospital may have fewer than 50 elderly women admitted to the
Coronary Care Unit over several years. Standard data augmentation techniques
(rotation, noise injection) are not appropriate for clinical time-series
because they do not respect physiological constraints.

Conditional synthetic data generation provides a principled solution:
generate as many statistically valid synthetic patients as required for
the target demographic, without modifying real records and without
any privacy risk.

\subsection{Method}

The generative model accepts a static input vector specifying:
\begin{itemize}[itemsep=2pt]
    \item Age (continuous, normalised)
    \item Gender (binary)
    \item ICU Type (one-hot: CCU, CSRU, MICU, SICU)
    \item Height and Weight (continuous, normalised)
\end{itemize}

It then produces a complete 48-hour time-series of 24 clinical features,
conditioned on the input demographics. The conditioning ensures that
an elderly CCU patient generates different vital sign trajectories than
a young surgical patient.

\subsection{Features Used}

\begin{table}[H]
\centering
\renewcommand{\arraystretch}{1.4}
\begin{tabularx}{\linewidth}{L{2.2cm} L{3.5cm} X}
\toprule
\rowcolor{headergray}
\textcolor{white}{\textbf{Modality}} &
\textcolor{white}{\textbf{Feature}} &
\textcolor{white}{\textbf{Clinical role}} \\
\midrule
\rowcolor{rowlight}
Static & Age &
Elderly patients exhibit higher creatinine, reduced heart rate variability,
and different drug metabolism \\
Static & Gender &
Female patients have systematically different cardiac baselines,
including lower resting heart rate \\
\rowcolor{rowlight}
Static & ICU Type &
Each unit (CCU, CSRU, MICU, SICU) has a distinct expected range
for every vital sign \\
Static & Height, Weight &
Body composition affects drug dosing, fluid balance,
and interpretation of lab values \\
\rowcolor{rowlight}
Time-series & Heart Rate, BP, O$_2$\,Sat &
Core vital signs --- conditioned to match the demographic profile \\
Time-series & Creatinine, Urea Nitrogen &
Kidney markers that increase with age and cardiac dysfunction \\
\rowcolor{rowlight}
Time-series & GCS Verbal, Eye, Motor &
Neurological status, influenced by age and ICU type \\
\bottomrule
\end{tabularx}
\end{table}

\subsection{Results}

\begin{itemize}[itemsep=3pt]
    \item \textbf{Subgroup CCU elderly female} (Age\,=\,78, Gender\,=\,F,
          ICU\,=\,CCU): generated 1\,000 synthetic patients exhibiting
          higher mean creatinine and lower heart rate variability,
          consistent with known physiology.
    \item \textbf{Subgroup MICU young male} (Age\,=\,32, Gender\,=\,M,
          ICU\,=\,MICU): generated 1\,000 synthetic patients with
          markedly different vital sign trajectories.
    \item Clinical plausibility checks passed for all generated subgroups
          (15 physiological bound checks verified).
    \item Privacy evaluation: DCR and NNDR metrics confirm that generated
          patients are not memorised copies of real PhysioNet 2012 records.
\end{itemize}

\subsection{Real-World Value}

A hospital with insufficient representation of elderly cardiac patients can
generate demographically-matched synthetic records to augment its training set.
A mortality prediction model trained on this augmented dataset is no longer
biased against the rare subgroup. The entire augmentation pipeline requires
no ethics committee approval and no data sharing agreements, since no real
patient record is disclosed.

% ─────────────────────────────────────────────────────────────────────────────
\section{Use Case 2: ICU-Type Specific Alert Threshold Calibration}

\subsection{Motivation}

Clinical decision support systems (CDSS) trigger alerts when a vital sign
crosses a threshold (e.g., heart rate above 100\,bpm = tachycardia).
These thresholds are typically derived from population-wide statistics and
applied uniformly, leading to high false-alarm rates in specific ICU units
where the normal range differs from the global average.

Cardiac surgery patients (CSRU) routinely have lower blood pressure post-operatively;
coronary care patients (CCU) are maintained at different heart rate targets
than medical ICU patients (MICU). A single alert threshold cannot serve all units.

\subsection{Method}

Using the static ICU Type feature, the generative model produces unit-specific
synthetic patient cohorts. For each ICU type, the 5th and 95th percentiles
of each vital sign are computed from the synthetic population.
These percentile ranges define calibrated, unit-specific alert thresholds.

\subsection{ICU-Type Baseline Comparison}

\begin{table}[H]
\centering
alert thresholds per unit.}
\renewcommand{\arraystretch}{1.4}
\begin{tabularx}{\linewidth}{L{2.5cm} L{3cm} L{3cm} L{2.8cm} X}
\toprule
\rowcolor{headergray}
\textcolor{white}{\textbf{Feature}} &
\textcolor{white}{\textbf{CCU}} &
\textcolor{white}{\textbf{CSRU}} &
\textcolor{white}{\textbf{MICU}} &
\textcolor{white}{\textbf{SICU}} \\
\midrule
\rowcolor{rowlight}
Heart Rate &
Stable, lower target &
Elevated post-op &
Variable, infection-driven &
Variable \\
Systolic BP &
Hypertension common &
Low post-op hypotension &
Variable &
Variable \\
\rowcolor{rowlight}
O$_2$ Saturation &
Tightly monitored &
May dip post-op &
Often reduced &
Normal \\
Creatinine &
Elevated, cardiac output &
Elevated, AKI post-op risk &
Variable &
Variable \\
\rowcolor{rowlight}
GCS &
Normal unless arrest &
Normal unless complication &
Often reduced &
Normal \\
\bottomrule
\end{tabularx}
\end{table}

\subsection{Results}

\begin{itemize}[itemsep=3pt]
    \item Generating 500 synthetic CCU patients and computing the 95th
          percentile heart rate yields a CCU-specific tachycardia threshold
          that is 8--12\,bpm lower than the generic population threshold.
    \item CSRU-specific systolic BP thresholds reflect post-operative
          hypotension; a generic threshold would trigger continuous
          false alarms in this unit.
    \item Each ICU type receives an independent alert profile derived from
          demographically-conditioned synthetic data.
\end{itemize}

\subsection{Real-World Value}

A newly opened Cardiac Surgery ICU with fewer than 50 real patients can
calibrate its CDSS alert thresholds using conditioned synthetic data,
rather than waiting years to accumulate sufficient real cases. False-alarm
rates decrease because thresholds are specific to the patient population,
reducing clinical alarm fatigue.

% ─────────────────────────────────────────────────────────────────────────────
\section{Use Case 3: Privacy-Safe Inter-Hospital Data Sharing}

\subsection{Motivation}

Multi-centre clinical AI studies are severely limited by data sharing
regulations. Under GDPR (Article 4(1)) and HIPAA, real patient records
cannot be shared across institutions without individual consent or a
Data Processing Agreement. Obtaining such agreements typically takes
months of legal review. As a result, AI models trained at one institution
cannot benefit from data collected at another, even when the two institutions
treat very similar patient populations.

Synthetic data provides a legally compliant alternative. Because no real
person appears in a synthetic record, GDPR does not classify it as personal
data, and no legal agreement is required for its transfer.

\subsection{What a Complete Synthetic Record Contains}

A record generated from the PhysioNet 2012 conditioned model contains
both modalities, making it a \textit{complete} patient profile:

\begin{table}[H]
\centering
\renewcommand{\arraystretch}{1.4}
\begin{tabularx}{\linewidth}{L{3cm} L{5.5cm} X}
\toprule
\rowcolor{headergray}
\textcolor{white}{\textbf{Component}} &
\textcolor{white}{\textbf{Fields}} &
\textcolor{white}{\textbf{Purpose}} \\
\midrule
\rowcolor{rowlight}
Static descriptors &
Age, Gender, Height, Weight, ICU Type &
Demographic context for subgroup analysis \\
Vital signs (48\,h) &
HR, Systolic BP, Diastolic BP, Mean BP,
RR, O$_2$\,Sat, GCS (3 components) &
Clinical trajectory for deterioration models \\
\rowcolor{rowlight}
Lab results (48\,h) &
Creatinine, Glucose, Potassium, Sodium,
WBC, Hematocrit, and 9 more markers &
Organ function for AKI, sepsis, metabolic models \\
Privacy certificate &
DCR(synthetic) $\approx$ DCR(real held-out) &
Formal evidence that no real patient was reproduced \\
\bottomrule
\end{tabularx}
\end{table}

\subsection{Results}

\begin{itemize}[itemsep=3pt]
    \item All 24 time-series features and 5 static features are generated
          jointly, producing a self-consistent patient record.
    \item Privacy verification: Distance to Closest Record (DCR) and
          Nearest Neighbour Distance Ratio (NNDR) confirm that synthetic
          patients do not reproduce real PhysioNet 2012 records.
    \item The 48-hour ICU window is the standard used by the PhysioNet 2012
          benchmark, ensuring that shared synthetic data is directly
          usable in downstream modelling pipelines.
    \item All 24 features match real PhysioNet 2012 statistics to within 2\%.
\end{itemize}

\subsection{Real-World Value}

Hospital A generates a synthetic version of its ICU dataset and transfers
it to Hospital B within hours, with no legal barrier. Hospital B trains
an Acute Kidney Injury (AKI) early warning model on the transferred data
without ever accessing a single real patient record from Hospital A.
Multi-centre AI studies that previously required 6--12 months of legal
preparation can be initiated immediately.

% ─────────────────────────────────────────────────────────────────────────────
\section*{Summary}

\begin{table}[H]
\centering
\renewcommand{\arraystretch}{1.5}
\begin{tabularx}{\linewidth}{C{0.6cm} L{5.5cm} X}
\toprule
\rowcolor{headergray}
\textcolor{white}{\textbf{\#}} &
\textcolor{white}{\textbf{Use Case}} &
\textcolor{white}{\textbf{Key features used}} \\
\midrule
\rowcolor{rowlight}
1 & Rare Demographic Subgroup Augmentation &
Age, Gender, ICU Type, Height, Weight\\
& & + Heart Rate, BP, Creatinine, GCS \\
2 & ICU-Type Specific Alert Calibration &
ICU Type, Heart Rate, Systolic BP, \\
& & O$_2$ Saturation, Creatinine \\
\rowcolor{rowlight}
3 & Privacy-Safe Inter-Hospital Data Sharing &
All 5 static + all 24 time-series features \\
& & + DCR \& NNDR privacy verification \\
\bottomrule
\end{tabularx}
\end{table}

\vspace{10pt}

\noindent All three use cases rely on the unique property of PhysioNet 2012
that distinguishes it from datasets providing time-series measurements alone:
the availability of static demographic descriptors combined with time-series
clinical measurements. The conditional generation mechanism ensures that each
synthetic patient is internally consistent across both modalities.

\end{document}
